\chapter{Complexity — Path Hard Run}
%%% generated by the blueprint pipeline from TCSlib.Complexity.SpaceComplexity.Logspace.PathHardRun
\section{Overview}

Declarations of the Lean module \texttt{TCSlib.Complexity.SpaceComplexity.Logspace.PathHardRun}.

\section{Declarations}

\begin{theorem}[Parsing and run behavior]
\lean{Complexity.PathHard.parse_run}
\label{Complexity.PathHard.parse_run}
\leanok
\uses{Complexity.LogProg.AHalt, Complexity.LogProg.AReach, Complexity.LogProg.AReach.halt, Complexity.LogProg.AReach.refl, Complexity.LogProg.AReach.trans, Complexity.LogProg.GuardLe, Complexity.LogProg.alignedPair_loop, Complexity.LogProg.inputLength_loop, Complexity.PathHard.A, Complexity.PathHard.St.Bd, Complexity.PathHard.St.toV, Complexity.PathHard.St.zero, Complexity.PathHard.advF, Complexity.PathHard.ans_halt, Complexity.PathHard.orc, Complexity.PathHard.orc_bit, Complexity.PathHard.orc_len, Complexity.PathHard.rP, Complexity.PathHard.rY, Turing.pairDecode, Turing.pairDecode_pairEncode, Turing.pairEncode}
\difficulty{4}
Let $a,b,f,B$ be natural numbers, let $m$ be a Boolean mode, let $N$ be a finite
nondeterministic Turing machine over the Boolean alphabet, and let $y$ be a finite
Boolean string with length at most $B$. If the pair parser fails on $y$, the program
starting in its initial instruction with all registers zero and no recorded answer halts
with answer false, and every configuration before halting satisfies the guard that its
instruction meets the program precondition and every register is at most $B$. Moreover,
for every pair of finite Boolean strings $x,\beta$ whose pairing is $y$, the same
initial run reaches the copy instruction assigning the input length to the value
register, with the length register equal to $|y|$, the scan pointer equal to $2|x|+2$,
all other registers zero, and no answer recorded; every earlier configuration satisfies
the same guard.
\end{theorem}

\begin{definition}[Initial machine configuration]
\lean{Complexity.PathHard.setSt}
\label{Complexity.PathHard.setSt}
\leanok
\uses{Complexity.PathHard.St}
Given natural numbers $Y,P,X,s,E$, define a machine configuration whose fields, in
order, are $y=Y$, $p=P$, $x=X$, $t=0$, $c=0$, $s=s$, $e=E$, $n=2\cdot 2^E$,
$m=2(2^E)^2$, $k=4(2^E)^2$, $q=P$, $i=0$, $d=0$, $u=0$, and $z=0$.
\end{definition}

\begin{theorem}[Bounded setup run]
\lean{Complexity.PathHard.setup_run}
\label{Complexity.PathHard.setup_run}
\leanok
\uses{Complexity.Codec.halvCount_loop, Complexity.Codec.incChain, Complexity.LogProg.AReach, Complexity.LogProg.AReach.trans, Complexity.LogProg.GuardLe, Complexity.LogProg.reach_of_eq, Complexity.LogProg.registerScale_loop, Complexity.PathHard.A, Complexity.PathHard.St.Bd, Complexity.PathHard.St.Bd.toV, Complexity.PathHard.St.toV, Complexity.PathHard.St.toV_c, Complexity.PathHard.St.toV_t, Complexity.PathHard.St.zero, Complexity.PathHard.cpDst, Complexity.PathHard.cpNext, Complexity.PathHard.cpSrc, Complexity.PathHard.cp_run, Complexity.PathHard.eInc, Complexity.PathHard.ell, Complexity.PathHard.orc, Complexity.PathHard.rC, Complexity.PathHard.rE, Complexity.PathHard.rS, Complexity.PathHard.rT, Complexity.PathHard.sEinc, Complexity.PathHard.sSinc, Complexity.PathHard.setSt, Complexity.PathHard.shExit, Complexity.PathHard.shReg, Complexity.PathHard.sh_run, Complexity.PathHard.step_of, Complexity.PathHard.two_pow_ell_le}
\difficulty{5}
Let $a,b,f,X,Y,B$ be natural numbers with $a\le b$, $2X+2\le Y$, and
$Y+8\bigl(2^{b\ell(X)+f}\bigr)^2\le B$. For any Boolean mode, finite nondeterministic
Turing machine $N$, and Boolean word $y$, the register program with its oracle
determined by $N$ has a finite run from the copy phase, starting with $y$-register $Y$,
pointer $2X+2$, and all other registers zero, to the width-check phase. Its final
registers are $Y$, $2X+2$, $X$, $0$, $0$, $a\ell(X)$, $b\ell(X)+f$,
$2\cdot2^{b\ell(X)+f}$, $2(2^{b\ell(X)+f})^2$, $4(2^{b\ell(X)+f})^2$, $2X+2$, $0$, $0$,
$0$, $0$, in register order; at every configuration before the endpoint, the program's
guard holds and every register is at most $B$.
\end{theorem}

\begin{theorem}[Decoding run dichotomy]
\lean{Complexity.PathHard.check_decode_run}
\label{Complexity.PathHard.check_decode_run}
\leanok
\uses{Complexity.Codec.bitsVal_le_of_width, Complexity.Codec.canon_of_getLastD, Complexity.Codec.getLastD_of_canon, Complexity.Codec.pe_getD_snd, Complexity.LogProg.AHalt, Complexity.LogProg.AReach, Complexity.LogProg.AReach.halt, Complexity.LogProg.AReach.trans, Complexity.LogProg.Canon, Complexity.LogProg.GuardLe, Complexity.LogProg.bitsVal, Complexity.LogProg.boundedDecode_loop, Complexity.LogProg.canonicalTail_loop, Complexity.LogProg.reach_of_eq, Complexity.PathHard.A, Complexity.PathHard.Lb, Complexity.PathHard.St, Complexity.PathHard.St.Bd, Complexity.PathHard.St.Bd.of_toV, Complexity.PathHard.St.Bd.toV, Complexity.PathHard.St.eq_qk, Complexity.PathHard.St.eq_yid, Complexity.PathHard.St.toV, Complexity.PathHard.St.toV_k, Complexity.PathHard.St.toV_p, Complexity.PathHard.St.toV_q, Complexity.PathHard.St.toV_y, Complexity.PathHard.ans_halt, Complexity.PathHard.dInc, Complexity.PathHard.orc, Complexity.PathHard.orc_bit, Complexity.PathHard.rD, Complexity.PathHard.rI, Complexity.PathHard.rK, Complexity.PathHard.rP, Complexity.PathHard.rQ, Complexity.PathHard.rY, Complexity.PathHard.tAdv, Turing.length_pairEncode, Turing.pairEncode}
\difficulty{6}
Let $a,b,f,B\in\mathbb{N}$, let $m$ be a Boolean mode, let $N$ be a finite
nondeterministic Turing machine over the Boolean alphabet, and let $x,\beta,y$ be
Boolean words. Suppose $y$ is the pair encoding of $x$ and $\beta$, and let $st$ be a
register record whose fifteen entries are at most $B$. Assume that the query-length
register of $st$ equals $|y|$, its pair-scan pointer is $2|x|+2$, its width-check
pointer equals that pair-scan pointer, its index and decoding scratch registers are
zero, and its width cap $k$ satisfies $1\le k$ and $2k\le B$. Then either $|\beta|\le
|\mathrm{bits}(k)|$ and $\beta$ is canonical, in which case there exist $q',k'\le B$
such that the machine $A(a,b,f,m)$ with the oracle determined by $N$ reaches the label
for the $2n\le i$ comparison from the state labeled by the true width-check test, with
register values $st$ and no result recorded, while satisfying its guard at every earlier
time; the final register values are those of $st$ with the query-length register set to
the pair-scan pointer, the index set to the binary value of $\beta$, and the width-check
pointer and cap set to $q'$ and $k'$. Or the conjunction of that length bound and
canonicity fails, and the same initial configuration halts with result false while
satisfying its guard at every earlier time.
\end{theorem}
